\documentclass[12pt,a4paper]{article}
\usepackage[utf8]{inputenc}
\usepackage[russian]{babel}
\usepackage{amsmath, amssymb}
\usepackage{geometry}
\usepackage{booktabs}
\usepackage{array}
\usepackage{graphicx}
\usepackage{float}
\usepackage{hyperref}

\geometry{top=2cm, bottom=2cm, left=2.5cm, right=2.5cm}

\title{Оценка коэффициента полувыведения препарата I-131\\
\large Анализ экспоненциальных моделей распада}
\author{}
\date{}

\begin{document}

\maketitle

\section*{Введение}

Проведён анализ данных 103 пациентов, получивших радиоактивный йод I-131. 
Для каждого пациента измерена мощность эквивалентной дозы (МАЭД) в 4--7 временных точках 
(от 1 до 68 часов после введения). Цель — оценка коэффициента полувыведения ($T_{1/2}$) 
и подбор оптимальных математических моделей для аппроксимации кривой распада.

\section{Математические модели}

Рассмотрены три модели экспоненциального распада, описывающие зависимость МАЭД от времени $t$:

\subsection{Модель M1: Однокомпонентный экспоненциальный распад}

Наиболее простая модель, предполагающая один процесс выведения:

\begin{equation}
A(t) = A_0 \cdot e^{-\lambda t}
\end{equation}

где:
\begin{itemize}
    \item $A(t)$ — значение МАЭД в момент времени $t$, мкЗв/ч
    \item $A_0$ — начальная активность (интерцепт), мкЗв/ч
    \item $\lambda$ — константа скорости выведения, ч$^{-1}$
    \item $t$ — время после введения, ч
\end{itemize}

Коэффициент полувыведения:

\begin{equation}
T_{1/2} = \frac{\ln 2}{\lambda} \approx \frac{0.693}{\lambda}
\end{equation}

\subsection{Модель M2: Двухкомпонентный экспоненциальный распад}

Учитывает два процесса выведения (быстрый и медленный):

\begin{equation}
A(t) = A_1 \cdot e^{-\lambda_1 t} + A_2 \cdot e^{-\lambda_2 t}
\end{equation}

где:
\begin{itemize}
    \item $A_1, A_2$ — амплитуды быстрой и медленной компонент
    \item $\lambda_1, \lambda_2$ — константы скорости ($\lambda_1 > \lambda_2$)
\end{itemize}

Эффективный коэффициент полувыведения определяется как время, за которое 
модельная кривая падает до половины начального значения:

\begin{equation}
T_{1/2}^{\text{eff}} = \min\{t : A(t) = A(0)/2\}
\end{equation}

\subsection{Модель M3: Трёхкомпонентный экспоненциальный распад}

Наиболее сложная модель с тремя процессами выведения:

\begin{equation}
A(t) = A_1 \cdot e^{-\lambda_1 t} + A_2 \cdot e^{-\lambda_2 t} + A_3 \cdot e^{-\lambda_3 t}
\end{equation}

где $\lambda_1 > \lambda_2 > \lambda_3 > 0$.

\section{Методология}

\subsection{Подбор параметров}

Параметры моделей определялись методом нелинейного наименьших квадратов 
(алгоритм Levenberg--Marquardt из \texttt{scipy.optimize.curve\_fit}) с ограничениями:
\begin{itemize}
    \item Все амплитуды $A_i > 0$
    \item Все константы скорости $\lambda_i > 0$
    \item $\lambda_i \leq 5$ ч$^{-1}$ (ограничение сверху для устойчивости)
\end{itemize}

Начальные параметры подбирались через логарифмическую линейную регрессию.

\subsection{Критерии сравнения моделей}

Для отбора лучшей модели использовался информационный критерий Акаике (AIC):

\begin{equation}
\text{AIC} = n \cdot \ln\left(\frac{\text{RSS}}{n}\right) + 2k
\end{equation}

где:
\begin{itemize}
    \item $n$ — количество точек данных
    \item $\text{RSS} = \sum (y_i - \hat{y}_i)^2$ — сумма квадратов остатков
    \item $k$ — количество параметров модели
\end{itemize}

Дополнительно вычислялись $R^2$ (коэффициент детерминации) и RMSE (среднеквадратичная ошибка).

\section{Результаты}

\subsection{Статистика выборки}

\begin{itemize}
    \item Обработано пациентов: \textbf{103} из 204 строк (50\%)
    \item Пропущено: 101 строка (неполные данные, отсутствие столбцов МАЭД)
    \item Временные точки: 4--7 на пациента (мода: 6 точек)
\end{itemize}

\subsection{Распределение лучших моделей}

\begin{table}[H]
\centering
\begin{tabular}{lcc}
\toprule
\textbf{Модель} & \textbf{Лучшая для} & \textbf{Доля} \\
\midrule
M1 (однокомпон.) & 42 пациента & 40.8\% \\
M2 (двухкомпон.) & 44 пациента & 42.7\% \\
M3 (трёхкомпон.) & 17 пациентов & 16.5\% \\
\bottomrule
\end{tabular}
\caption{Распределение лучших моделей по критерию AIC}
\end{table}

\subsection{Коэффициент полувыведения ($T_{1/2}$)}

\begin{table}[H]
\centering
\begin{tabular}{lc}
\toprule
\textbf{Статистика} & \textbf{Значение} \\
\midrule
Минимум & 4.63 ч \\
Медиана & 13.61 ч \\
Среднее & 14.73 ч \\
Максимум & 43.08 ч \\
\bottomrule
\end{tabular}
\caption{Статистика коэффициента полувыведения I-131}
\end{table}

\subsection{Качество аппроксимации}

\begin{table}[H]
\centering
\begin{tabular}{lccc}
\toprule
\textbf{Метрика} & \textbf{Мин} & \textbf{Среднее} & \textbf{Макс} \\
\midrule
$R^2$ (лучшая модель) & 0.8406 & 0.9927 & 1.0000 \\
RMSE (лучшая модель) & 0.0166 & 0.4821 & 2.1534 \\
\bottomrule
\end{tabular}
\caption{Качество аппроксимации данных моделями}
\end{table}

\section{Резюме}

\subsection*{Тенденции}

\begin{enumerate}
    \item \textbf{Доминирование M2.} Двухкомпонентная модель признана лучшей для 42.7\% пациентов, 
    что подтверждает наличие двух фаз выведения I-131 из организма.
    
    \item \textbf{Высокое качество подгонки.} Среднее значение $R^2 = 0.9927$ свидетельствует 
    о том, что экспоненциальные модели адекватно описывают кривые распада.
    
    \item \textbf{Широкий разброс $T_{1/2}$.} Коэффициент полувыведения варьируется от 4.6 до 43.1 часа — 
    в 9.3 раза. Это указывает на значительную межличностную вариабельность метаболизма I-131.
    
    \item \textbf{Редкость M3.} Трёхкомпонентная модель лучше только для 16.5\% пациентов, 
    что говорит о том, что третья фаза выведения присутствует не у всех.
\end{enumerate}

\subsection*{Лучший вариант}

\textbf{M2 (двухкомпонентная экспонента)} — оптимальный баланс между точностью и простотой:
\begin{itemize}
    \item Лучшая модель для 44 пациентов (42.7\%)
    \item Учитывает две фазы выведения (быстрый — из крови/внеклеточной жидкости, 
    медленный — из щитовидной железы/тканей)
    \item Не переобучается (6 параметров против 4 у M2)
    \item Физически интерпретируема
\end{itemize}

\subsection*{Мнение}

\begin{itemize}
    \item Данные I-131 хорошо описываются биэкспоненциальной моделью, что согласуется 
    с физиологией: быстрый компонент отражает выведение из крови и мягких тканей, 
    медленный — из щитовидной железы.
    
    \item Разброс $T_{1/2}$ в 9 раз требует индивидуального подхода к дозировке. 
    Стандартные протоколы могут быть неоптимальны для пациентов с экстремальными значениями.
    
    \item Для 17 пациентов M3 даёт лучшее описание — возможно, это пациенты с 
    аномальным метаболизмом или дополнительными факторами накопления.
    
    \item Рекомендация: использовать M2 как основную модель, с переходом на M3 
    при $R^2_{\text{M2}} < 0.95$.
\end{itemize}

\end{document}